\bookchapter{第二章：整理你的數字房間——檔案管理}

\chapterintro{手機用久了會亂，好習慣從管理檔案開始。就像你的房間需要定期整理一樣，你的數字空間也需要打理。一個井井有條的手機，不僅找東西方便，執行也更流暢。}

\booksection{一、三層儲存體系：你的數字倉庫}

\begin{enumerate}
  \item \textbf{機身儲存}：核心區，速度快，系統和App都在這。這是你手機最寶貴的空間，要精打細算。
  \item \textbf{SD卡（如果支援）}：擴充衣櫃，速度慢點，適合放電影、照片等大檔案。注意：現在很多旗艦機不支援SD卡了。
  \item \textbf{雲端儲存}：保險櫃+同步網盤。防丟失，也能在多個裝置間同步。iCloud、百度網盤、OneDrive都是常用選擇。
\end{enumerate}

\booksection{二、認識“檔案管理器”}

安卓一般有個 [檔案] 或 [檔案管理] App，iOS是自帶的 [檔案] App。它們是你檢視手機所有檔案的“資源管理器”。

\booksubsection{常見資料夾含義：}

\begin{itemize}
  \item \textbf{DCIM}：相機拍攝的照片和影片。
  \item \textbf{Pictures}：其他圖片，如截圖、儲存的圖片。
  \item \textbf{Download}：瀏覽器下載的檔案。
  \item \textbf{Music}：音樂檔案。
  \item \textbf{Movies}：影片檔案。
  \item \textbf{Documents}：文件檔案。
  \item \textbf{Android}：安卓系統資料夾，裡面是各App的資料，一般不要手動刪除。
\end{itemize}

\booksection{三、檔案系統結構：深入瞭解}

安卓基於Linux，檔案系統有一定的層級結構。雖然普通使用者不需要深入瞭解，但知道一些基本概念有幫助。

\booksubsection{儲存分割區：}

\begin{itemize}
  \item \textbf{/system}：系統分割區，存放作業系統檔案。普通使用者無權訪問，需要root許可權。
  \item \textbf{/data}：資料分割區，存放使用者資料和應用資料。這是你日常使用的空間。
  \item \textbf{/sdcard 或 /storage/emulated/0}：使用者可訪問的儲存空間，你的照片、下載檔案都在這裡。
\end{itemize}

\booksubsection{App資料儲存：}

\begin{itemize}
  \item \textbf{內部儲存}：App的私有空間，其他App無法訪問。解除安裝App時會被刪除。
  \item \textbf{外部儲存}：App在使用者可訪問區域儲存的檔案，如照片、下載內容。解除安裝App時不一定會被刪除。
\end{itemize}

\booksection{四、資料備份與恢復：給你的資料上保險}

資料無價，備份是必須養成的習慣。

\booksubsection{備份策略：}

\begin{enumerate}
  \item \textbf{雲備份}：自動、方便，但需要網路和雲空間。iCloud、華為雲、小米雲等都支援。
  \item \textbf{本地備份}：用資料線連線電腦，用官方工具備份。速度快、空間大，但需要手動操作。
  \item \textbf{雙重備份}：雲備份+本地備份，最安全。
\end{enumerate}

\booksubsection{備份什麼？}

\begin{itemize}
  \item \textbf{照片和影片}：最佔空間，也最珍貴。
  \item \textbf{通訊錄}：存在雲賬號裡最方便。
  \item \textbf{微信聊天記錄}：用微信自帶的備份功能。
  \item \textbf{重要文件}：存在雲盤或電腦上。
\end{itemize}

\booksection{五、檔案安全：保護你的隱私}

\begin{itemize}
  \item \textbf{敏感檔案加密}：很多手機有"安全資料夾"或"私密空間"功能，可以把敏感檔案放進去，需要密碼或指紋才能訪問。
  \item \textbf{刪除不等於銷燬}：普通刪除只是標記空間可覆蓋，資料還在。要徹底刪除，需要用專門的安全擦除工具，或者多次覆蓋寫入。
  \item \textbf{雲端安全}：選擇信譽好的雲服務商，開啟雙重認證，定期檢查登入裝置。
\end{itemize}

\booksection{六、空間不夠？來場大掃除}

\begin{itemize}
  \item \textbf{清理快取}：App的臨時記憶，清了能騰地方。在設定 > 儲存 裡可以清理。
  \item \textbf{查詢大檔案}：用檔案管理器按大小排序，揪出佔空間的"元兇"。
  \item \textbf{刪除殘留}：解除安裝App後，有時會留下垃圾檔案，手動清掉。
  \item \textbf{清理相簿}：刪除重複照片、模糊照片、不需要的截圖。
  \item \textbf{清理微信}：微信是空間大戶，定期清理聊天記錄和快取很有必要。
\end{itemize}

\booksection{七、常見檔案格式全解析}

手機裡各種檔案字尾名讓人頭大，搞懂它們能幫你更好地管理檔案：

\begin{center}
\begin{tabularx}{\linewidth}{|X|X|X|X|}
\hline
\textbf{格式} & \textbf{型別} & \textbf{特性} & \textbf{推薦開啟方式} \tabularnewline
\hline
JPG/JPEG & 圖片 & 有失真壓縮，體積小，相容性最好。手機拍照預設格式。 & 相簿/任意圖片檢視器 \tabularnewline
PNG & 圖片 & 無失真壓縮，支援透明背景。截圖、圖示常用。體積比JPG大。 & 相簿/任意圖片檢視器 \tabularnewline
MP3 & 音訊 & 有失真壓縮，體積小，相容性極好。最通用的音訊格式。 & 任意音樂播放器 \tabularnewline
FLAC & 音訊 & 無失真壓縮，音質與CD一致，體積比MP3大3-5倍。發燒友首選。 & 音樂播放器/專業播放App \tabularnewline
WAV & 音訊 & 無壓縮原始音訊，音質最好但體積巨大。專業錄音常用。 & 音樂播放器 \tabularnewline
M4A & 音訊 & AAC編碼，比MP3同位元速率音質更好。蘋果生態常用。 & Apple Music/音樂播放器 \tabularnewline
MP4 & 影片 & 最通用的影片封裝格式，相容性極好。手機錄製的影片預設格式。 & 相簿/任意影片播放器 \tabularnewline
MKV & 影片 & 開源封裝格式，支援多音軌、多字幕。電影資源常用。 & VLC/nPlayer \tabularnewline
DOC/DOCX & 文件 & Word文件。DOCX是壓縮格式，體積更小。DOC是舊版格式。 & WPS/Office \tabularnewline
PPT/PPTX & 演示 & PowerPoint簡報。PPTX是壓縮格式。 & WPS/Office \tabularnewline
XLS/XLSX & 表格 & Excel電子表格。XLSX是壓縮格式。 & WPS/Office \tabularnewline
PDF & 文件 & 跨平臺格式，排版固定不可編輯。合同、論文、電子書常用。 & Adobe Reader/WPS \tabularnewline
ZIP & 壓縮 & 最通用的壓縮格式，跨平臺相容好。 & 系統自帶/解壓App \tabularnewline
RAR & 壓縮 & 壓縮率比ZIP高，但需要專用軟體解壓。 & 解壓App \tabularnewline
ISO & 映象 & 光碟映象檔案，完整複製光碟內容。系統安裝包常用。 & 解壓App/虛擬光碟機 \tabularnewline
HTML & 網頁 & 網頁檔案，用瀏覽器開啟。儲存的網頁常用此格式。 & 瀏覽器 \tabularnewline
\hline
\end{tabularx}
\end{center}

\begin{tcbpractice}{練一練：整理"下載"資料夾}
開啟檔案管理器 > 進入“下載”資料夾。看看裡面都有啥？把沒用的刪掉，有用的分門別類（比如新建“工作資料”、“表情包”資料夾放進去）。順便看看儲存空間，瞭解哪些型別檔案佔用最多。

\end{tcbpractice}

\booksection{八、圖片與檔案系統管理：讓資產井井有條}

上一節我們學會了“認識資料夾”，但光認識還不夠。本節要講的是\textbf{如何主動建立一套可長期執行的管理規範}，讓你的檔案隨時找得到、不丟失、不被偷窺。

\booksubsection{檔案分類與命名規範}

混亂的根源是沒有規則。建立規則其實很簡單，關鍵在於\textbf{堅持}。

\begin{itemize}
  \item \textbf{"時間+事件"命名法}：資料夾名字裡帶時間戳和事件描述，方便日後回憶與檢索。例如\texttt{202606\_旅行票據}、\texttt{2025年終總結}、\texttt{2024孩子生日}。這樣既能看到時間線，也能一眼知道內容。
  \item \textbf{三級分類體系}：
  \begin{itemize}
    \item \textbf{一級}：按檔案型別（文件 / 圖片 / 影片 / 音訊）
    \item \textbf{二級}：按使用場景（工作 / 生活 / 學習）
    \item \textbf{三級}：按具體事件或專案（如"2026Q2\_專案A"）
  \end{itemize}
  \item \textbf{月度清理計劃}：每月固定一個日期（建議發工資那天或月初1號），對"下載"目錄進行集中整理：\textbf{歸檔有用的、刪除無用的、轉存珍貴的}。一年 12 次小整理，比攢一年大掃除輕鬆得多。
\end{itemize}

\begin{tcbpractice}{本月任務清單}
挑一個空閒的週末，按以下順序操作：

\begin{enumerate}
  \item 在檔案管理器中建立\texttt{2026Q2\_專案A}形式的本月工作資料夾；
  \item 把所有"下載"目錄裡本週的檔案過一遍，留下有用的、歸檔、刪除沒用的；
  \item 給今年的家庭照片建立\texttt{2025家庭}、\texttt{2026家庭}這種年度資料夾。
\end{enumerate}

\end{tcbpractice}

\booksubsection{冗餘檔案處理策略}

手機用久了會“長胖”，多半是冗餘檔案在作怪。

\begin{itemize}
  \item \textbf{72 小時自動清理}：微信、QQ 等社交 App 接收的臨時檔案（表情包、聊天圖片、小影片），大多看一次就沒用了。在 [微信] > [設定] > [通用] > [儲存空間] 裡開啟"自動清理"，把過期快取設成 72 小時。手機瞬間清爽不少，還不丟重要記錄。
  \item \textbf{儲存空間定向清理}：安卓 [設定] > [儲存] > [清理加速]，iOS [設定] > [通用] > [iPhone 儲存空間]，系統會按"快取 / 臨時檔案 / 可解除安裝 App / 大檔案"分類列出。重點清理：(1) 快取檔案；(2) 已過期安裝包（升級 App 後留下的舊版安裝包）；(3) 重複檔案（同名同尺寸的照片）。
  \item \textbf{自動分類規則}：多數品牌相簿 App 支援智慧分類：人物、風景、文件、截圖。開啟後，新拍的合影會自動歸入"人物 > 張三"，掃描的合同會自動歸入"文件"。\textbf{配合每月整理，可極大降低"找不到東西"的機率}。
\end{itemize}

\booksubsection{圖片安全處理規範}

圖片是手機裡最“敏感”的資料型別——一張隨手拍就能暴露你的位置、裝置、時間。我們建議\textbf{建立“雙重保護”工作流}：

\begin{tcbwarn}{傳送證件照/敏感截圖的標準流程}
\begin{enumerate}
  \item \textbf{可見資訊打碼}：用系統自帶或第三方工具（如美圖秀秀、Snapseed）\textbf{塗抹身份證號中間幾位、車牌號、門牌、姓名、金額}等關鍵欄位；
  \item \textbf{後設資料剝離}：傳送前在相簿中點"分享"或"匯出"，系統會詢問是否保留位置資訊，選擇\textbf{"不保留"}。這樣拍攝裝置型號、GPS 座標、拍攝時間就不會跟著圖片一起發出去；
  \item \textbf{臨時水印}（可選）：在敏感截圖上加"僅供XX辦理使用"的文字水印，避免被濫用。
\end{enumerate}

\end{tcbwarn}

後設資料（EXIF）清理的常用方法（基礎操作，深度防護見第六章“照片脫敏處理”）：

\begin{itemize}
  \item \textbf{系統自帶}：iOS [設定] > [隱私與安全性] > [相機]，可關閉"精確位置"；分享時選"僅當前照片"會\textbf{自動剝離大部分 EXIF}。
  \item \textbf{安卓}：多數品牌相簿 App 在"分享"時會提示"是否移除位置資訊"，勾選即可。
  \item \textbf{專業工具}：ExifTool（電腦端）、Metapho（iOS）、Photo Metadata Remover（安卓）可以徹底清除所有後設資料。
\end{itemize}

對於需要\textbf{長期儲存}的敏感圖片（身份證、銀行卡、戶口本掃描件、護照），建議存入\textbf{加密相簿}：

\begin{itemize}
  \item iPhone：[照片] > 選中照片 > 隱藏（僅 Face ID 驗證後可檢視）；
  \item 華為/小米：[相簿] > 長按照片 > "移入保險櫃"；
  \item 第三方加密 App：\textbf{360 安全相簿}、\textbf{騰訊相簿管家（隱私相簿）}、\textbf{KeepSafe}等，單獨密碼保護，隱蔽性更好。
\end{itemize}
